\section{Problem Statement}
\label{sec:problem}
Current PBS (via MEV-Boost) improves proposer's equality but introduces new trust. Specifically:
\begin{itemize}
    \item \textbf{Relay Trust Centralization:} A few relay servers act as gatekeepers between builders and proposers. Builders must trust relays not to censor blocks or steal payloads; proposers must trust relays to provide the promised block. This concentration of trust subverts decentralization.
    \item \textbf{Builder-Proposer Trust:} After a proposer signs a builder’s header, the builder is supposed to deliver the full block. A malicious builder could withhold or delay the block (griefing the proposer and wasting its slot). Conversely, if a proposer could see the block before signing, it could steal the MEV by not publishing or by re-proposing an empty slot. MEV-Boost’s design disallows the former, but a cryptographically enforced solution is desirable.
    \item \textbf{Economic Guarantee:} In MEV-Boost, a builder commits to pay a bid to the proposer, but proving that the block \emph{indeed} pays that bid (e.g. via included fee transactions or a bonded deposit) requires trusting that the builder will honor it. Ideally, the system should ensure the bid is actually covered by the block’s contents in a verifiable way.
    \item \textbf{Centralizing Incentives:} As shown by, heterogeneity in block production can lead to concentration among top builders, even under PBS. If only the largest builders reliably win auctions, power can centralize in those entities, such entities can delay and censor transactions according to their desire. A fully trustless PBS should not worsen this trend and might even mitigate it by lowering entry barriers (as more and more builders will be able to join the protocol and build blocks, by lowering the trust entry barrier).
\end{itemize}
To summarize, a trustless PBS would enable proposers to (a) obtain assurance that the block they are purchasing is valid and finalizable, and (b) verify that the corresponding payment will be executed as promised. Our approach addresses these challenges through optimistic zero-knowledge proofs, which replace the need for trusted intermediaries such as relays and builders with cryptographic verification.
